\begin{lemma}
For every prefix \(p\) and every \(r\ge0\), if
\(\ell_r^p\) is the layer choice of \(\noderef{StarProductLayerChoice}\),
then
\[
  |U_r^p|\le (r+1)|Z_{\ell_r^p}^p|.
\]
\end{lemma}
\begin{proof}
By \(\noderef{StarProductRankAtMostSet}\), the set \(U_r^p\) consists of
the points \(y\) whose rank \(r^p(y)\) is at most \(r\).  By
\(\noderef{StarProductRankLayer}\), the layer \(Z_s^p\) consists of the
points whose rank is exactly \(s\).  Hence every \(y\in U_r^p\) belongs to
the layer \(Z_{r^p(y)}^p\), and the index \(r^p(y)\) lies in
\(\{0,\ldots,r\}\).  Thus \(U_r^p\) is contained in the union of the
\(r+1\) layers \(Z_s^p\) for \(0\le s\le r\).

The cardinality of a finite union is at most the sum of the cardinalities of
the sets being unioned.  Therefore
\[
  |U_r^p|\le \sum_{s=0}^r |Z_s^p|.
\]
For each \(s\le r\), \(\noderef{StarProductLayerChoiceMax}\) gives
\(|Z_s^p|\le |Z_{\ell_r^p}^p|\).  The sum has \(r+1\) terms, so it is at
most \((r+1)|Z_{\ell_r^p}^p|\).  Combining the two inequalities proves the
claim.
\end{proof}
